% !TeX root = ../../Book3.tex
% 纲要标题：### H.4 椭圆曲线的几何
% 内容：亏格一曲线的方程、由除子类定义的群律及同源的基本计算。

\section{椭圆曲线的几何}
\label{b3:sec:H04}

亏格为一时，正次数除子的截面数恰好等于次数。因此，允许函数在一个指定点有二阶、三阶极点，便能选出两个函数；它们之间最早出现的关系给出三次方程。另一方面，一次除子类的唯一点代表又把除子类的加法变成曲线上点的运算。方程与群律由此来自同一组 Riemann--Roch 计算。

\ProseParagraphBreak
本节仍以代数闭域 \(k\) 上的光滑射影整曲线为对象。由除子类定义点的群运算不限制特征；短 Weierstrass 方程、以下坐标公式及同源计算均假设 \(\operatorname{char}k\ne2,3\)。方程构造与群结构可对照\cite[Chapter IV, Proposition 4.6, Proposition 4.8, pp. 319--322]{Hartshorne1977}。

\subsection{Weierstrass 方程与光滑平面三次曲线}
\label{b3:subsec:H04:01}

\begin{definition}\label{b3:def:H-elliptic-curve}
	设 \(k\) 为代数闭域，\(E\) 为光滑射影整曲线，\(O\in E(k)\) 为一个指定点。若 \(g(E)=1\)，则称 \((E,O)\) 为一条\term[tuoyuan-quxian]{椭圆曲线}[elliptic curve]。指定点 \(O\) 称为原点。椭圆曲线的同构要求是把指定原点送到指定原点的 \(k\) 曲线同构。
\end{definition}

这里的亏格是 \(h^1(E,\mathcal O_E)\)，并不是方程中某个变量的次数。原点也不是额外的一组坐标；它将确定函数可以出现极点的位置，以及后面的群单位元。

\begin{theorem}[Weierstrass 模型]\label{b3:thm:H-weierstrass-model}
	设 \(k\) 为代数闭域，\(\operatorname{char}k\ne2,3\)，\((E,O)\) 为 \(k\) 上的椭圆曲线。则存在 \(A,B\in k\)，使 \((E,O)\) 同构于平面曲线
	\begin{equation}\label{b3:eq:H-short-weierstrass}
		Y^2Z=X^3+AXZ^2+BZ^3,\qquad O=[0:1:0],
	\end{equation}
	其中
	\[
		\Delta\coloneqq-16(4A^3+27B^2)\ne0.
	\]
	反过来，满足这个条件的平面三次方程与指定点 \([0:1:0]\) 构成一条椭圆曲线。
\end{theorem}
\begin{proof}
	由\cref{b3:cor:H-canonical-degree}，\(E\) 上的典范除子满足 \(K\sim0\)。对 \(n\ge1\)，\(K-nO\) 的次数为负，故没有非零截面；\cref{b3:thm:H-riemann-roch}给出
	\[
		\ell(nO)=n.
	\]
	特别地，\(L(O)=k\)。选取
	\[
		x\in L(2O)\setminus L(O),\qquad
		y\in L(3O)\setminus L(2O).
	\]
	它们在 \(O\) 的极点阶数分别恰为二、三，在其余点均正则。于是 \(1,x,y,x^2,xy,x^3\) 的极点阶数依次为 \(0,2,3,4,5,6\)，彼此不同，因而线性无关。它们与 \(y^2\) 都属于六维空间 \(L(6O)\)，所以有一条线性关系。关系中 \(y^2\) 与 \(x^3\) 的系数均非零，否则剩余项的最高极点阶数不能相消。适当缩放 \(y\) 后，可写成
	\[
		y^2+a_1xy+a_3y=x^3+a_2x^2+a_4x+a_6.
	\]
	代数闭性保证所需的非零标量有平方根。因二在 \(k\) 中可逆，把 \(y\) 换成 \(y+(a_1x+a_3)/2\)，可消去左端的一次 \(y\) 项；再因三可逆，平移 \(x\) 可消去右端的二次项。这些变换不改变 \(x,y\) 的极点阶数，所得关系即为 \(y^2=x^3+Ax+B\)。

	还须说明所得方程确实表示原来的光滑曲线。\Cref{b3:prop:H-high-degree-linear-systems}适用于次数为三的除子 \(3O\)，给出由完整线性系 \(|3O|\) 定义的闭嵌入
	\[
		E\longrightarrow\mathbb P^2_k,
		\qquad P\longmapsto[x(P):y(P):1]\quad(P\ne O).
	\]
	在 \(O\) 附近用局部统一化元 \(t\) 同乘 \(t^3\)，三个截面变为 \(t^3x,t^3y,t^3\)，前后两项在 \(O\) 消失，中间一项为单位，所以 \(O\) 的像为 \([0:1:0]\)。这里闭嵌入所用的输入是 \(3O\) 分离不同点及同一点的一阶切向，而不仅是三个函数线性无关。三次多项式 \(x^3+Ax+B\) 在射影直线的无穷远处有奇数阶极点，故在 \(k(x)\) 中不是平方。因此 \(y^2-x^3-Ax-B\) 不可约；闭嵌入的像遂恰为\eqref{b3:eq:H-short-weierstrass}的整个三次曲线。

	设 \(f(x)=x^3+Ax+B\)。无穷远处只有 \(O\)，且齐次方程对 \(Z\) 的偏导数在 \(O\) 为一。仿射部分出现奇点，当且仅当
	\[
		y=0,\qquad f(x)=f'(x)=0.
	\]
	三次多项式的判别式为 \(-4A^3-27B^2\)，故光滑性等价于 \(\Delta\ne0\)。也可直接核验：重根 \(r\) 必满足 \(A=-3r^2\)、\(B=2r^3\)；反向在 \(A\ne0\) 时取 \(r=-3B/(2A)\)，在 \(A=0\) 时取 \(r=0\)。原来的曲线光滑，因而它的方程必满足这一条件。

	最后证明反向的亏格断言。对满足 \(\Delta\ne0\) 的方程，有理微分
	\[
		\omega=\frac{\mathrm{d}x}{2y}
	\]
	在仿射开集 \(y\ne0\) 上是局部微分模的生成元；在 \(y=0\) 处，由 \(f'(x)\ne0\) 及 \(2y\,\mathrm{d}y=f'(x)\,\mathrm{d}x\)，它又写成 \(\mathrm{d}y/f'(x)\)，仍是生成元。在 \(O\) 附近取 \(u=X/Y\)、\(v=Z/Y\)，方程为
	\[
		v=u^3+Auv^2+Bv^3.
	\]
	其中 \(u\) 是统一化元，\(1-2Auv-3Bv^2\) 为单位。微分上式并代入 \(x=u/v\)、\(y=1/v\)，得到
	\[
		\omega=-\frac{\mathrm{d}u}{1-2Auv-3Bv^2}.
	\]
	所以 \(\omega\) 处处无零点、无极点，其除子为零。再由\cref{b3:cor:H-canonical-degree}的 \(\deg K=2g-2\)，得到 \(g=1\)。
\end{proof}

上述方程称为\term[duan-Weierstrass-fangcheng]{短 Weierstrass 方程}[short Weierstrass equation]，\(\Delta\) 称为这一方程的判别式。证明中二与三的可逆性分别用于配方和平移消去二次项；特征二、三的椭圆曲线仍有 Weierstrass 方程，但不能统一采用这个短式。沿用证明中的 \(u,v\)，还有 \(\operatorname{ord}_O(v)=3\)，从而 \(\operatorname{ord}_O(x)=-2\)、\(\operatorname{ord}_O(y)=-3\)，与最初的截面选择一致。

\begin{example}\label{b3:exm:H-weierstrass-divisors}
	设 \(k\) 为代数闭域，\(\operatorname{char}k\ne2,3\)，考虑 \(E:y^2=x^3-x\)。这里 \(\Delta=64\)，故曲线光滑。记 \(T_a=(a,0)\)，其中 \(a=-1,0,1\)。在每个 \(T_a\) 处，\(y\) 是统一化元，\(x-a\) 的阶数为二；在 \(O\) 处，\(x,y\) 的极点阶数为二、三。因此
	\[
		\operatorname{div}(x)=2T_0-2O,
		\qquad
		\operatorname{div}(y)=T_{-1}+T_0+T_1-3O.
	\]
	这些是函数除子的等式。下一小节将把它们转写为点的群运算关系。
\end{example}

\subsection{割线、切线与群律}
\label{b3:subsec:H04:02}

一条直线与光滑三次曲线相交，计重数得到三个点。用这条直线作函数除子，所得关系属于除子类群；这才是割线与群运算之间的联系。为区别点的运算和除子的形式相加，暂把点的运算记作 \(P\oplus Q\)。

\begin{theorem}[椭圆曲线的群律]\label{b3:thm:H-elliptic-group-law}
	设 \(k\) 为代数闭域，\((E,O)\) 为 \(k\) 上的椭圆曲线。映射
	\[
		\iota:E(k)\longrightarrow\operatorname{Pic}^0(E),
		\qquad P\longmapsto[P-O]
	\]
	为双射。存在唯一以 \(O\) 为单位元的交换群运算 \(\oplus\)，使 \(\iota\) 为群同构；具体地，
	\[
		P\oplus Q=R\quad\Longleftrightarrow\quad P+Q\sim R+O.
	\]
	若另有 \(\operatorname{char}k\ne2,3\)，并采用\eqref{b3:eq:H-short-weierstrass}，则任一直线的交点除子 \(P+Q+R\) 满足
	\[
		P\oplus Q\oplus R=O,
	\]
	其中允许重复点，并按交点重数计数。
\end{theorem}
\begin{proof}
	对任意次数零除子类 \([D]\)，\(D+O\) 的次数为一。由\cref{b3:prop:H-genus-one-sections}，它的类中存在唯一点 \(P\)，于是 \([D]=[P-O]\)。这同时证明 \(\iota\) 的满射性和单射性。把 \(\operatorname{Pic}^0(E)\) 的加法经此双射传回 \(E(k)\)，便得到所述运算。因为目标是交换群，结合律、交换律与逆元的存在均随之成立；\([O-O]=0\) 给出单位元。定义等式
	\[
		[P-O]+[Q-O]=[R-O]
	\]
	又正好等价于 \(P+Q\sim R+O\)。

	在上述三次嵌入中，直线 \(Z=0\) 截出的除子为 \(3O\)。任意另一条直线 \(\ell=0\) 与它的比 \(\ell/Z\) 是曲线上的有理函数，其主除子为 \(D_\ell-3O\)。所以 \(D_\ell\sim3O\)，且 \(\deg D_\ell=3\)。写作 \(D_\ell=P+Q+R\) 后，便有 \([P-O]+[Q-O]+[R-O]=0\)。这里用的是截面零除子的次数，不需要把实平面上的图形当作结合律的证明。
\end{proof}

对于仿射点 \(P=(a,b)\)，竖直直线 \(X-aZ=0\) 截出 \(P+(a,-b)+O\)，计重数的说法也包括 \(b=0\)。因此
\[
	-P=(a,-b),\qquad -O=O.
\]
两点确定的割线若在第三点 \(R\) 相交，则 \(P\oplus Q=-R\)；两点重合时用切线，并把这个点计至少两次。特别地，\(O\) 的切线 \(Z=0\) 与曲线的交点除子为 \(3O\)，所以 \(O\) 是拐点。

\begin{proposition}[割线与切线公式]\label{b3:prop:H-chord-tangent-formulas}
	设 \(k\) 为代数闭域，\(\operatorname{char}k\ne2,3\)，\(E:y^2=x^3+Ax+B\) 满足 \(\Delta\ne0\)，原点为 \(O\)。对仿射点 \(P=(x_1,y_1)\)、\(Q=(x_2,y_2)\)，若 \(x_1\ne x_2\)，令
	\[
		m=\frac{y_2-y_1}{x_2-x_1};
	\]
	若 \(P=Q\) 且 \(y_1\ne0\)，令
	\[
		m=\frac{3x_1^2+A}{2y_1}.
	\]
	在这两种情形下，\(P\oplus Q=(x_3,y_3)\)，其中
	\begin{equation}\label{b3:eq:H-addition-formulas}
		x_3=m^2-x_1-x_2,\qquad
		y_3=m(x_1-x_3)-y_1.
	\end{equation}
	其余情形由 \(P\oplus O=P\) 及 \(P\oplus(-P)=O\) 决定；特别地，\(y_1=0\) 时 \(P\oplus P=O\)。
\end{proposition}
\begin{proof}
	非竖直的割线或切线写成 \(y=m(x-x_1)+y_1\)。代入曲线方程，得到首一三次方程
	\[
		x^3-m^2x^2+\text{一次及常数项}=0.
	\]
	它的三个根是 \(x_1,x_2,x_R\)，切线时前两个根相同。因此 \(x_R=m^2-x_1-x_2\)。第三交点的纵坐标为 \(m(x_R-x_1)+y_1\)，取逆元即得 \(y_3\)。切线斜率由 \(2y_1\,\mathrm{d}y=(3x_1^2+A)\,\mathrm{d}x\) 给出。若 \(x_1=x_2\)，方程使 \(y_2=y_1\) 或 \(y_2=-y_1\)；除已列出的非竖直切线情形外，均落入竖直直线的情形，故已穷尽。
\end{proof}

\begin{example}\label{b3:exm:H-point-addition}
	在 \(k=\overline{\mathbb Q}\) 上取 \(E:y^2=x^3-x+1\)，有 \(\Delta=-16\cdot23\ne0\)。点 \(P=(0,1)\)、\(Q=(1,1)\) 之间的割线为 \(y=1\)，第三交点为 \((-1,1)\)，所以 \(P\oplus Q=(-1,-1)\)。在 \(P\) 处的切线斜率为 \(-1/2\)，从而
	\[
		[2]P\coloneqq P\oplus P
		=\left(\frac14,-\frac78\right).
	\]
	代回右端可核验 \((1/4)^3-1/4+1=49/64\)。两种计算分别使用割线和切线，所得运算的结合律已经由除子类群保证。
\end{example}

点集上有群运算，还不足以说明它与代数几何相容。特别是\eqref{b3:eq:H-addition-formulas}的分母消失时，不能把分式直接当成已定义的坐标。

\begin{proposition}\label{b3:prop:H-elliptic-group-morphisms}
	设 \(k\) 为代数闭域，\(\operatorname{char}k\ne2,3\)，\((E,O)\) 为 \(k\) 上的椭圆曲线。取逆映射 \(E\rightarrow E\)、加法映射 \(\mu:E\times E\rightarrow E\) 以及每个平移 \(\tau_A:P\mapsto A\oplus P\) 均为 \(k\) 态射；平移是曲线自同构。
\end{proposition}
\begin{proof}
	先固定 \(A\in E(k)\)。由\cref{b3:prop:H-genus-one-sections}和\cref{b3:prop:H-high-degree-linear-systems}，次数二的除子 \(A+O\) 有两个线性无关截面，且无基点，因而给出态射 \(f_A:E\rightarrow\mathbb P^1_k\)，纤维除子属于 \(|A+O|\)，次数为二。由\cref{b3:prop:H-curve-morphism}，它是二次有限态射。由于特征不为二，函数域扩张是可分二次扩张，有唯一非平凡自同构。应用\cref{b3:lem:H-rational-map-extension}于对应的有理自映射及其逆，得到 \(E\) 的自同构 \(\sigma_A\)：两个复合在函数域上为恒等，由延拓唯一性处处为恒等。

	在不分歧纤维中，\(\sigma_A\) 交换两个点；在分歧纤维中，它固定唯一的点，纤维为该点的二倍。因此对每个 \(P\) 都有
	\[
		P+\sigma_A(P)\sim A+O,
		\qquad \sigma_A(P)=A\oplus(-P).
	\]
	于是 \(\sigma_O\) 就是取逆映射，而 \(\tau_A=\sigma_A\circ\sigma_O\)。这证明了取逆和平移的代数性。

	在 \(E\times E\) 的非空开集 \(U\) 上，要求两个点均为仿射点且横坐标不同；\eqref{b3:eq:H-addition-formulas}在那里定义态射 \(\mu_U:U\rightarrow E\)。对任意闭点 \((P,Q)\)，可选 \(A,B\in E(k)\)，使 \((P\oplus A,Q\oplus B)\in U\)。例如先选 \((P_0,Q_0)\in U\)，再令 \(A=P_0\oplus(-P)\)、\(B=Q_0\oplus(-Q)\)。于是
	\[
		\tau_{-(A\oplus B)}\circ\mu_U\circ(\tau_A\times\tau_B)
	\]
	在包含 \((P,Q)\) 的开邻域上是态射，并在每个点处都等于抽象群运算 \(\oplus\)。这些邻域覆盖整个 \(E\times E\)：其补集若非空，就是有限型 \(k\) 概形中的非空闭集，必有闭点，与构造矛盾。这些局部态射在重叠处相同，遂黏合成 \(\mu\)。这里相同也可按态射核验：源为约化 \(k\) 概形，目标分离，闭点上的相等迫使态射相等。其结合律和单位元等恒等式同样由闭点上的群公理得到。
\end{proof}

这个证明用平移把非空开集上的加法公式移到任意一对点的邻域，补上了分母消失处的定义。二次线性系和这一延拓方法见\cite[Chapter IV, Lemma 4.2, Proposition 4.8, pp. 318, 321--322]{Hartshorne1977}。

\subsection{j 不变量与同源}
\label{b3:subsec:H04:03}

短 Weierstrass 方程并不唯一。改变 \(L(2O)\)、\(L(3O)\) 的基会改变系数；因此，要比较两条曲线，须先确定哪些系数变化只是在改写同一条曲线。

\begin{proposition}[短方程的同构与 \(j\) 不变量]\label{b3:prop:H-j-invariant}
	设 \(k\) 为代数闭域，\(\operatorname{char}k\ne2,3\)，两条椭圆曲线分别为
	\[
		E:y^2=x^3+Ax+B,\qquad
		E':y'^2=x'^3+A'x'+B',
	\]
	两式判别式均非零，原点均为无穷远点。每个保持原点的同构 \(E\rightarrow E'\) 恰有形式
	\[
		(x,y)\longmapsto(u^2x,u^3y),\qquad
		A'=u^4A,\quad B'=u^6B,
		\quad u\in k^\times.
	\]
	数值
	\begin{equation}\label{b3:eq:H-j-invariant}
		j(E)\coloneqq1728\,\frac{4A^3}{4A^3+27B^2}
	\end{equation}
	只依赖椭圆曲线的同构类，称为\term[j-bubianliang]{\(j\) 不变量}[j-invariant]。两条这样的椭圆曲线同构，当且仅当 \(j(E)=j(E')\)。
\end{proposition}
\begin{proof}
	同构保持 \(O\) 处的极点阶数，故它在函数上的拉回必有形式
	\[
		x'=ax+r,\qquad y'=by+sx+t,\qquad a,b\ne0.
	\]
	把它代入 \(E'\) 的方程，并用 \(y^2=x^3+Ax+B\) 约化。因为 \(1,y\) 在 \(k(x)\) 上线性无关，比较一次 \(y\) 项得 \(s=t=0\)；比较 \(x^2\) 项得 \(r=0\)；比较 \(x^3\) 项得 \(b^2=a^3\)。令 \(u=b/a\)，便有 \(u^2=a\)、\(u^3=b\)。比较其余两项即得 \(A'=u^4A\)、\(B'=u^6B\)。反向的缩放及其逆显然给出保持原点的同构。

	缩放使 \(A^3,B^2\) 同乘 \(u^{12}\)，所以\eqref{b3:eq:H-j-invariant}不变。现在设两条曲线的 \(j\) 相同。若 \(j=0\)，则 \(A=A'=0\)，而 \(B,B'\ne0\)，取 \(u^6=B'/B\) 即可。若 \(j=1728\)，则 \(B=B'=0\)，而 \(A,A'\ne0\)，取 \(u^4=A'/A\) 即可。其余情形四个系数均非零。置 \(r=A'/A\)、\(s=B'/B\)，相同的 \(j\) 给出 \(r^3=s^2\)。取 \(a=s/r\)，则 \(a^2=r\)、\(a^3=s\)；再取 \(u^2=a\)，便同时满足 \(u^4=r\)、\(u^6=s\)。所需根的存在用到了 \(k\) 的代数闭性。
\end{proof}

例如，\(y^2=x^3-x\) 的 \(j\) 为 \(1728\)，\(y^2=x^3+1\) 的 \(j\) 为零。若只比较不带原点的曲线同构，一样得到这个分类：给定同构 \(f:E\rightarrow E'\)，再复合目标上的平移 \(\tau_{-f(O)}\)，就得到保持原点的同构。这里 \(j\) 的分类范围是代数闭域上的光滑曲线同构；它不是同源的分类。Hartshorne 采用 Legendre 方程给出同一不变量的分类\cite[Chapter IV, Theorem 4.1, pp. 317--320]{Hartshorne1977}。

\begin{definition}\label{b3:def:H-isogeny}
	设 \(k\) 为代数闭域，\(\operatorname{char}k\ne2,3\)，\((E,O)\)、\((E',O')\) 为 \(k\) 上的椭圆曲线。一个非恒定 \(k\) 态射 \(\varphi:E\rightarrow E'\)，若满足
	\[
		\varphi(O)=O',\qquad
		\varphi(P\oplus Q)=\varphi(P)\oplus\varphi(Q)
		\quad(P,Q\in E(k)),
	\]
	则称为一个\term[tongyuan]{同源}[isogeny]。其次数定义为
	\[
		\deg\varphi=[k(E):\varphi^*k(E')],
	\]
	包括函数域扩张的不可分次数。若该扩张可分，则称 \(\varphi\) 为可分同源。记 \(\ker\varphi(k)=\{\,P\in E(k)\mid\varphi(P)=O'\,\}\)，这是有限的点群；概形纤维 \(\varphi^{-1}(O')\) 还保留各点的重数。
\end{definition}

在这个定义中，群条件其实由保持原点的非恒定曲线态射自动满足。下面给出所需的除子理由，而不把它作为额外约定。

\begin{lemma}\label{b3:lem:H-origin-morphism-homomorphism}
	设 \(k\) 为代数闭域，\(\operatorname{char}k\ne2,3\)，\((E,O)\)、\((E',O')\) 为椭圆曲线。每个满足 \(f(O)=O'\) 的非恒定 \(k\) 态射 \(f:E\rightarrow E'\) 都是同源。
\end{lemma}
\begin{proof}
	由\cref{b3:prop:H-curve-morphism}，\(f\) 有限且平坦。将点除子 \(P\) 推为 \(f(P)\)，并线性延拓，得到 \(f_*:\operatorname{Div}(E)\rightarrow\operatorname{Div}(E')\)。这个推送保持线性等价，因为对 \(h\in k(E)^\times\)，有
	\[
		f_*\operatorname{div}(h)
		=\operatorname{div}\bigl(\operatorname{Norm}_{k(E)/f^*k(E')}(h)\bigr).
	\]
	说明这个等式的局部依据：在目标点的 DVR 上，有限平坦代数是有限秩自由模，范数是乘 \(h\) 的线性映射的行列式。对局部正则的非零 \(h\)，在 DVR 上逐次选取赋值最小的非零矩阵项作主元，以可逆行列变换消去同行同列的其余项，便把矩阵化为对角形。对角元阶数之和同时等于行列式的赋值与余核长度；按源中的各点分解，这一长度是各 \(\operatorname{ord}_P(h)\) 之和，因剩余域均为 \(k\)。对一般有理函数清去分母并相减，即得所写等式。这也适用于不可分扩张。

	若 \(P\oplus Q=R\)，则 \(P+Q\sim R+O\)。推送后得到
	\[
		f(P)+f(Q)\sim f(R)+O',
	\]
	由\cref{b3:thm:H-elliptic-group-law}，这恰好是 \(f(P)\oplus f(Q)=f(R)\)。
\end{proof}

这一结论也见\cite[Chapter IV, Lemma 4.9, p. 322]{Hartshorne1977}。同源的纤维是核的平移，但一个纤维的点数和它的次数仍可能不同。

\begin{proposition}[同源的次数与核]\label{b3:prop:H-isogeny-degree}
	设 \(k\) 为代数闭域，\(\operatorname{char}k\ne2,3\)，\(\varphi:(E,O)\rightarrow(E',O')\) 为同源，次数为 \(d\)。则 \(\varphi\) 有限、平坦且满射。设 \(e\) 为它在 \(O\) 处的分歧指数，则在每个闭点的分歧指数都为 \(e\)，并且
	\[
		d=e\,\#\ker\varphi(k).
	\]
	下列条件等价：\(\varphi\) 可分；\(e=1\)；\(\varphi\) 在每个点都不分歧；\(\#\ker\varphi(k)=d\)。
\end{proposition}
\begin{proof}
	前三项性质由\cref{b3:prop:H-curve-morphism}给出。对 \(P\in E(k)\)，群同态性质给出交换图
	\[
		\begin{tikzcd}
			E \arrow[r,"\tau_P"] \arrow[d,"\varphi"'] & E \arrow[d,"\varphi"] \\
			E' \arrow[r,"\tau_{\varphi(P)}"'] & E' .
		\end{tikzcd}
	\]
	水平箭头均为自同构，所以 \(\varphi\) 在 \(P\) 与 \(O\) 处的局部参数阶数相同。这说明各闭点分歧指数相同。对原点纤维应用\cref{b3:prop:H-curve-morphism}的纤维次数公式，得 \(d=\sum_{P\in\ker\varphi(k)}e=e\,\#\ker\varphi(k)\)。

	令 \(L=\varphi^*k(E')\)、\(K=k(E)\)。两条曲线光滑，故 \(K\otimes_L\Omega_{L/k}\) 与 \(\Omega_{K/k}\) 都是一维 \(K\) 空间。\Cref{b3:thm:differential-transitivity}给出正合列
	\[
		K\otimes_L\Omega_{L/k}\longrightarrow\Omega_{K/k}
		\longrightarrow\Omega_{K/L}\longrightarrow0.
	\]
	第一条映射非零，恰好等价于 \(\Omega_{K/L}=0\)；由\cref{b3:thm:finite-field-differentials-separable}，这又等价于 \(K/L\) 可分。

	若 \(\varphi\) 可分，函数域处的微分非零，故在某个非空开集上的切映射非零。对局部统一化元，切映射非零恰好表示拉回有非零一次项，也就是分歧指数为一。平移使不分歧性传到全部点，所以 \(e=1\)。反过来，若 \(e=1\)，\(O\) 处的切映射非零，函数域处的微分便非零，由刚证的判据得到可分性。其余等价关系随已证公式得到。曲线态射的这一微分计算也见\cite[Chapter IV, Propositions 2.1--2.2, pp. 299--300]{Hartshorne1977}。
\end{proof}

\begin{example}[倍点同源]\label{b3:exm:H-doubling-isogeny}
	设 \(k\) 为代数闭域，\(\operatorname{char}k\ne2,3\)，\(E:y^2=x^3+Ax+B\) 满足 \(\Delta\ne0\)。反复复合加法和取逆，得到整数倍映射 \([n]:E\rightarrow E\)。它在点上是交换群同态。在 \((O,O)\) 处，加法的微分限制到两个坐标轴都是恒等，所以
	\[
		\mathrm{d}\mu_{(O,O)}(v,w)=v+w,
		\qquad \mathrm{d}[n]_O=n\,\operatorname{id}_{T_OE}.
	\]
	特别地，\([2]\) 的微分非零，故它非恒定且可分。又因 \([2]P=O\) 等价于 \(P=-P\)，其点核恰为
	\[
		E[2](k)=\{\,O,(r_1,0),(r_2,0),(r_3,0)\,\},
	\]
	其中 \(r_1,r_2,r_3\) 为 \(x^3+Ax+B\) 的三个不同根。于是\cref{b3:prop:H-isogeny-degree}给出 \(\deg[2]=4\)，而且 \(E[2](k)\cong(\mathbb Z/2\mathbb Z)^2\)。在 \(y\ne0\) 的仿射部分，\([2]\) 的坐标由\eqref{b3:eq:H-addition-formulas}取 \(P=Q\) 算出；在其余点的值由已经建立的态射与群律决定。
\end{example}

\symidx{\([n]\)：椭圆曲线的整数倍映射}

若正特征整除 \(n\)，上面的切映射为零，不能据此认定 \([n]\) 恒定，也不能直接用核的点数计算次数。在特征五的三次曲线上，可以直接算出一个不可分同源。

\begin{example}\label{b3:exm:H-inseparable-isogeny}
	取 \(k=\overline{\mathbb F}_5\)，\(E:Y^2Z=X^3-XZ^2\)，原点为 \(O=[0:1:0]\)。因方程系数属于 \(\mathbb F_5\)，坐标式
	\[
		F:E\longrightarrow E,
		\qquad [X:Y:Z]\longmapsto[X^5:Y^5:Z^5]
	\]
	定义一个保持原点的非恒定 \(k\) 态射，故由\cref{b3:lem:H-origin-morphism-homomorphism}，它是同源。它在 \(k\) 点上为双射，特别地 \(\ker F(k)=\{O\}\)。然而在 \(O\) 的统一化元 \(u=X/Y\) 上，\(F^*u=u^5\)，所以 \(e=5\)，进而 \(\deg F=5\)。所有坐标微分均为零，与不可分性相符。这个 \(k\) 态射在坐标环上固定 \(k\) 的元素；它不同于把常数也取五次方的绝对 Frobenius 环映射。
\end{example}

三次方程可以用来核验光滑性、计算群运算和比较同构类；同源则比较保持群结构的有限态射，其次数必须连同可分性一起判断。一般整数倍映射的次数、挠点的完整结构及椭圆曲线的算术理论需要进一步的工具，这里只使用已经计算的倍点与五次方态射。

% BEGIN APPENDIX OUTLINE SYNC
% * 状态：初稿已完成；新增范围尚未正式审核。
% #### H.4.1 Weierstrass 方程与光滑平面三次曲线
% #### H.4.2 割线、切线与群律
% #### H.4.3 j 不变量与同源
% ---
% END APPENDIX OUTLINE SYNC
